% Changed: restore service-object argument and verified serving evidence.
% Protected: qualification semantics, measurement clocks, adverse results and citation keys.
% Meaning-Check: PRESERVED
% Risk-Flags: none
\section{Conclusion}
\label{sec:conclusion}
ContractKV governs reusable KV as a service object whose layout, behavioral evidence and transactional visibility share one validity boundary. CALC exposes a dependency-valid physical prefix; scoped certification establishes admissibility; PSTP controls what becomes visible; and the Service-Template Certificate binds authorized reuse to a versioned lifecycle. This coupling gives the service a concrete decision about which transformed execution artifacts later tenants may consume.

The measurements establish complementary parts of that argument. Query-qualified execution achieves 1.070--1.619$\times$ full-cost throughput across three models, with unchanged IDs in all 1,296 admitted compiled future pairs. Deployment-qualified registration reduces calls 25-fold and unhidden verification exposure 7.7-fold; a separate open-arrival point gains 0.658 compliant requests/s, while 80-identity density replay uses 71.9\% fewer peak KV blocks. Qualification therefore has a reusable economic unit as well as a publication boundary. Three lessons follow: admission scope determines shareability, atomic commit governs visibility, and locality plus residency determine how much reuse is realized. Aligning these decisions is the service-level contribution of ContractKV.
